\datedsubsection{ Fev. 2025 -- Fev. 2028 }
	{%
		École nationale des Ponts et chaussées, ERC Discover, \textit{AIKON}
	}
	{%
		\textbf{Ingénieur de recherche}
	}
	{%
		Au sein de l'\href{https://erc-discover.github.io/}{ERC Discover}, j'ai participé au développement de la \href{https://aikon-platform.github.io/}{plateforme AIKON}, constituée de plusieurs applications Web intégrant \textit{computer vision} et gestion de corpus historiques. Pour traiter les volumes massifs de données produits par la vision artificielle, j'ai créé le serveur d'annotations IIIF \href{https://github.com/Aikon-platform/aiiinotate/tree/main}{aiiinotate}, conçu pour être très rapide, scalable à plusieur dizaines de millions d'annotations et compatible avec Mirador 4.
	}
\datedsubsection{ Nov. 2022 -- Déc. 2024 }
	{%
		Institut national d'histoire de l'art, \textit{Richelieu. Histoire du quartier}}
	{%
		\textbf{Ingénieur d'études}~}
	{%
		Ingénieur d'études principal du projet \href{https://quartier-richelieu.inha.fr/}{Rich.Data}, j'ai été chargé de la conception et de la mise en \oe{}uvre de la chaîne de traitement du projet: modélisation et administration d'une base PostgreSQL, développement \textit{full-stack} d'une plateforme Web d'exploration et analyse de données (API REST développée avec Flask, \textit{frontend} Vue.js), migration de données.}
\datedsubsection{ Avril 2022 -- Sept. 2022 }
	{%
		École normale supérieure, \textit{MSS/Katabase}}
	{%
		\textbf{Stage de fin d'études}~ École nationale des Chartes}
	{%
		Mise en \oe{}uvre et enrichissement de la chaîne de traitement du projet MSS/Katabase: gestion d'un corpus de 500 documents XML-TEI, fouille de texte et analyse quantitative, développement applicatif: refonte du \href{https://katabase.huma-num.fr/}{site}, mise au point d'une API REST.}
\datedsubsection{ Juin 2021 -- Sept. 2021 }
	{%
		Archives nationales, \textit{ANR DesignSHS}}
	{%
		\textbf{Stage}~}
	{%
		Stagiaire archiviste dans le cadre d’une ANR: modélisation et enrichissement d'une base
		de données, exploration de perspectives en analyse statistique et visualisation de données.}
\datedsubsection{ Mars 2021 -- Mai 2021 }
	{%
		INA GRM (Groupe de recherches musicales)}
	{%
		\textbf{Stage de fin d'études}~École du Louvre}
	{%
		Documentation des productions bibliographiques du GRM. Modélisation d'une base de données, étude comparative de systèmes de gestion de bases de données \textit{open-source}; réflexions sur l'établissement d’un thésaurus et l'archivage de logiciels de création musicale.}
\datedsubsection{ Sept. 2020 -- Déc. 2020 }
	{%
		Galerie Marian Goodman}
	{%
		\textbf{Stage}~}
	{%
		Asssistant en galerie: gestion de l'inventaire informatisé, mise en place de la stratégie de communication, relations presse.}